\chapter{Closed Cells in Sections B}
\label{ch:assignment3}

\begingroup\itshape Assignment 3 --- 19 points\endgroup

\workedby{\tbd{naam}}

\bigskip
\tbd{Inleiding van dit hoofdstuk: de drie stijfferconfiguraties van het
torsiekaster/dek, de afmetingen uit de opgave (T-stijffers, trapeziums,
plaatdiktes), de opsplitsing in open en gesloten delen en de aannames
($b\gg t_p$).}

\section{Shear flow sketch for Configurations 1 and 2}
\label{sec:a3-1}

\questionbox{Sketch the shear flow using arrows for Configurations 1 and 2.
\textbf{(2 pts)}}

$M_t$ is applied counter-clockwise, so the shear flow in the closed cells also runs counter-clockwise (Figure~\ref{fig:a3flow}). The T-stiffeners are open sections: they are only attached with their web, so there is no closed path for a shear flow. Like we saw in the lectures, the shear stress in an open section only goes up and down within the plate thickness, so the T-stiffeners carry no
shear flow $q$.

\paragraph{Configuration 1.}
The box is one closed cell, so the shear flow $q$ is constant along all four
plates.

\paragraph{Configuration 2.}
The trapezoids form three extra closed cells with the bottom plate. In the legs and top of the trapezoids the flows of the two cells go against each other, so only the small difference $q_m - q_t$ is left there.


\begin{figure}[h]
    \centering
    \begin{tikzpicture}[scale=2.0, >={Stealth[length=2mm,width=1.6mm]}]
        \tikzset{
          big/.style={->, red!75!black, line width=2.2pt},
          mid/.style={->, red!75!black, line width=1.8pt},
          small/.style={->, red!75!black, line width=0.7pt},
          stiff/.style={line width=0.8pt}}
        % T-stiffener: \Tdown{x}{y_plate}{height}, \Tup{x}{y_plate}{height}
        \def\Tdown#1#2#3{\draw[stiff] (#1,#2) -- (#1,{#2-#3});
                         \draw[stiff] ({#1-0.125},{#2-#3}) -- ({#1+0.125},{#2-#3});}
        \def\Tup#1#2#3{\draw[stiff] (#1,#2) -- (#1,{#2+#3});
                       \draw[stiff] ({#1-0.125},{#2+#3}) -- ({#1+0.125},{#2+#3});}

        % ======================= Configuration 1 =======================
        \begin{scope}[yshift=2.75cm]
            \draw[line width=1.2pt] (0,0) rectangle (5,2);
            \foreach \x in {0.833,1.667,2.5,3.333,4.167}{\Tdown{\x}{2}{0.5}\Tup{\x}{0}{0.5}}
            % shear flow (constant)
            \foreach \x in {1.25,2.917,4.583}{
                \draw[big] ({\x-0.2},0) -- ({\x+0.2},0);
                \draw[big] ({\x+0.2},2) -- ({\x-0.2},2);}
            \draw[big] (5,0.8) -- (5,1.2);
            \draw[big] (0,1.2) -- (0,0.8);
            \node[red!75!black, below] at (2.917,-0.05) {$q$};
            % torque
            \draw[->, thick] (2.5,1) ++(-30:0.3) arc (-30:210:0.3);
            \node at (2.5,1) {\small $M_t$};
            \node[anchor=east] at (-0.3,1) {(a)};
        \end{scope}

        % ======================= Configuration 2 =======================
        \begin{scope}
            \draw[line width=1.2pt] (0,0) rectangle (5,2);
            \foreach \x in {0.833,1.667,2.5,3.333,4.167}{\Tdown{\x}{2}{0.5}}
            % trapezoids: base from x0 to x0+5/6, height b/4, 45 deg legs
            \foreach \x in {0.625,2.083,3.542}{
                \draw[line width=0.6pt] (\x,0) -- ({\x+0.25},0.25) -- ({\x+0.5833},0.25) -- ({\x+0.8333},0);
                \Tup{\x+0.4167}{0.25}{0.25}
                % net flow in legs and crown (q_m - q_t), thin
                \draw[small] ({\x+0.07},0.07) -- ({\x+0.18},0.18);
                \draw[small] ({\x+0.30},0.25) -- ({\x+0.40},0.25);
                \draw[small] ({\x+0.763},0.07) ++(-0.11,0.11) -- ({\x+0.763},0.07);
                % base (q_t)
                \draw[mid] ({\x+0.3},0) -- ({\x+0.53},0);}
            % bottom plate between trapezoids (q_m)
            \foreach \x in {0.31,1.77,3.23,4.69}{\draw[big] ({\x-0.15},0) -- ({\x+0.15},0);}
            % top plate and sides (q_m)
            \foreach \x in {1.25,2.917,4.583}{\draw[big] ({\x+0.2},2) -- ({\x-0.2},2);}
            \draw[big] (5,0.8) -- (5,1.2);
            \draw[big] (0,1.2) -- (0,0.8);
            % labels
            \node[red!75!black, below] at (1.77,-0.05) {$q_m$};
            \node[red!75!black, below] at (2.5,-0.05) {$q_t$};
            \draw[very thin, gray] (2.62,0.27) -- (3.0,0.75);
            \node[red!75!black, above, font=\small] at (3.05,0.72) {$q_m - q_t$};
            % torque
            \draw[->, thick] (2.5,1.05) ++(-30:0.3) arc (-30:210:0.3);
            \node at (2.5,1.05) {\small $M_t$};
            \node[anchor=east] at (-0.3,1) {(b)};
        \end{scope}
    \end{tikzpicture}
    \caption{Shear flow for a counter-clockwise torque $M_t$ in (a) Configuration~1
    and (b) Configuration~2. The arrow thickness gives the size of the shear flow.
    The T-stiffeners are open sections and carry no shear flow.}
    \label{fig:a3flow}
\end{figure}

\finalanswer{See Figure~\ref{fig:a3flow}. Configuration~1 has one constant shear flow around the box. In Configuration~2 the box plating carries $q_m$, the bottom plate under the trapezoids $q_t$ and the trapezoid legs and top the small difference $q_m - q_t$. The open T-stiffeners carry no shear flow.}

\newpage
\section{T-stiffener versus L-stiffener}
\label{sec:a3-2}

\questionbox{Explain the consequences for the torsion stiffness contribution of open sections if a T-stiffener is exchanged for an L-stiffener, in case the web and flange dimensions remain the same. \textbf{(1 pt)}}

The T-stiffener is an open section, so its torsion stiffness follows from Saint-Venant for thin plates:
\begin{equation}
  I_{p,open} = \frac{1}{3}\sum_i b_i\,t_i^3 .
\end{equation}
This only depends on the length and thickness of the web and flange, not on how they are connected. In the lecturesslides from week 2 page 46. we saw that no matter how the plates are connected or oriented, $I_p$ stays the same. So an L-stiffener with the same web and flange gives exactly the same contribution to the torsion stiffness.

\finalanswer{There is no difference. $I_{p,open} = \frac{1}{3}\sum b_i t_i^3$ only
depends on the plate lengths and thicknesses, so the T- and L-stiffener give the
same torsion stiffness contribution.}

\section{Shear flow components}
\label{sec:a3-3}

\questionbox{Calculate the shear flow components $q_i(b,t_p,M_t)$ for
Configurations 1 and 2. \textbf{(4 pts)}}

\tbd{Uitwerking}

\finalanswer{\tbd{kort en expliciet antwoord}}

\section{Maximum shear stress}
\label{sec:a3-4}

\questionbox{\begin{enumerate}[label=\alph*., leftmargin=*]
  \item Locate the position of maximum shear stress and obtain
        $\tau_{max,i}(b,t_p,M_t)$ for Configurations 1 and 2 (assuming
        $b\gg t_p$). \textbf{(3 pts)}
  \item Compare the outcomes of Configurations 1 and 2 to the (given) outcome of
        Configuration 3 and discuss the differences. \textbf{(2 pts)}
\end{enumerate}}

\begin{enumerate}[label=\alph*., leftmargin=*]
  \item \tbd{Uitwerking}
  \item \tbd{Uitwerking}
\end{enumerate}

\finalanswer{\tbd{kort en expliciet antwoord per deelvraag}}

\section{Specific torsion stiffness and strength}
\label{sec:a3-5}

\questionbox{\begin{enumerate}[label=\alph*., leftmargin=*]
  \item Calculate the (specific) torsion stiffness $I_{p,i}(b,t_p)$ and strength
        $M_{t,i}(b,t_p,\tau_{max})$ for Configurations 1 and 2. \textbf{(3 pts)}
  \item Compare the outcomes of Configurations 1 and 2 to the (given) outcome of
        Configuration 3 and discuss the origin of possible differences.
        \textbf{(2 pts)}
  \item For Configurations 1 and 2 split the relative stiffness and strength
        contributions for the open and closed parts and discuss the differences.
        \textbf{(2 pts)}
\end{enumerate}}

\begin{enumerate}[label=\alph*., leftmargin=*]
  \item \tbd{Uitwerking}
  \item \tbd{Uitwerking}
  \item \tbd{Uitwerking}
\end{enumerate}

\finalanswer{\tbd{kort en expliciet antwoord per deelvraag}}
